% !TEX program = xelatex
% c5p-s-2026-final-answer.tex
% 2026 €€£ 超专业级别硬件能力认证第一轮（C5P-S1）提高级 C++语言试题·参考答案与解析
\documentclass[a4paper]{ctexart}

\usepackage{fontspec}
\usepackage{geometry}
\geometry{top=2.7cm, left=1.9cm, right=1.9cm, bottom=2.9cm}
\setmainfont{Consolas}
\setCJKmainfont[BoldFont=SimHei]{SimSun}
\usepackage{tikz}
\usetikzlibrary{arrows.meta,positioning}
\usepackage{eso-pic}
\newcommand{\waterword}{%
  \makebox[0pt][c]{%
    \tikz[baseline]{%
      \node[inner sep=0pt, rotate=-35, text=black, opacity=0.085,
        font=\fontsize{54}{70}\selectfont] {折工};%
    }%
  }%
}
\AddToShipoutPictureFG{%
  \put(\LenToUnit{-0.3cm},\LenToUnit{1.2cm}){\waterword}%
  \put(\LenToUnit{5.0cm},\LenToUnit{0.0cm}){\waterword}%
  \put(\LenToUnit{10.3cm},\LenToUnit{1.4cm}){\waterword}%
  \put(\LenToUnit{15.6cm},\LenToUnit{0.2cm}){\waterword}%
  \put(\LenToUnit{2.2cm},\LenToUnit{6.2cm}){\waterword}%
  \put(\LenToUnit{7.5cm},\LenToUnit{4.8cm}){\waterword}%
  \put(\LenToUnit{12.8cm},\LenToUnit{6.3cm}){\waterword}%
  \put(\LenToUnit{18.1cm},\LenToUnit{4.9cm}){\waterword}%
  \put(\LenToUnit{-0.3cm},\LenToUnit{10.0cm}){\waterword}%
  \put(\LenToUnit{5.0cm},\LenToUnit{11.5cm}){\waterword}%
  \put(\LenToUnit{10.3cm},\LenToUnit{10.1cm}){\waterword}%
  \put(\LenToUnit{15.6cm},\LenToUnit{11.6cm}){\waterword}%
  \put(\LenToUnit{2.2cm},\LenToUnit{16.5cm}){\waterword}%
  \put(\LenToUnit{7.5cm},\LenToUnit{15.1cm}){\waterword}%
  \put(\LenToUnit{12.8cm},\LenToUnit{16.6cm}){\waterword}%
  \put(\LenToUnit{18.1cm},\LenToUnit{15.2cm}){\waterword}%
}
\usepackage{fancyhdr}
\usepackage{lastpage}
\usepackage{amsmath}
\usepackage{booktabs}
\usepackage{longtable}
\usepackage{xcolor}
\DeclareMathSizes{10.53937}{10.95}{7.5}{5.5}
\setlength{\parindent}{0pt}
\setlength{\parskip}{3pt}
\linespread{1.25}

\newcommand{\ans}[2]{\noindent\textbf{#1} \quad 答案：#2\par}
\newcommand{\expn}[1]{\hspace*{1.5em}解析：#1\par}
\newcommand{\bigtitle}[1]{\begin{center}\dengxian\large #1\end{center}}
\newcommand{\overview}[1]{\par\smallskip\noindent\textbf{程序总析：}#1\par\smallskip}
\newCJKfontfamily{\dengxian}{DengXian}

% ---------------- 阅读程序配套图表 ----------------
\newcommand{\kmpchart}{%
\begin{center}
\begin{tikzpicture}[
  prebox/.style={draw=blue!55!black, fill=blue!5, rounded corners=2pt,
    minimum width=1.35cm, minimum height=0.52cm, font=\scriptsize},
  chain/.style={draw=orange!70!black, fill=orange!8, rounded corners=2pt,
    minimum width=1.05cm, minimum height=0.48cm, font=\small},
  flow/.style={-{Latex[length=2mm]}, thick, draw=black!65}
]
\node[anchor=west, font=\small\bfseries] at (0,1.08) {前缀函数示例：\texttt{ababa}};
\node[prebox] (p1) at (0.70,0.35) {\texttt{a}\;|\;\(\pi=0\)};
\node[prebox] (p2) at (2.25,0.35) {\texttt{ab}\;|\;\(\pi=0\)};
\node[prebox] (p3) at (3.80,0.35) {\texttt{aba}\;|\;\(\pi=1\)};
\node[prebox] (p4) at (5.35,0.35) {\texttt{abab}\;|\;\(\pi=2\)};
\node[prebox] (p5) at (6.90,0.35) {\texttt{ababa}\;|\;\(\pi=3\)};
\draw[flow] (p1) -- (p2);
\draw[flow] (p2) -- (p3);
\draw[flow] (p3) -- (p4);
\draw[flow] (p4) -- (p5);
\node[anchor=west, font=\small] at (0,-0.35) {回退链：};
\node[chain] (c5) at (1.25,-0.35) {长度 5};
\node[chain] (c3) at (3.00,-0.35) {长度 3};
\node[chain] (c1) at (4.75,-0.35) {长度 1};
\node[chain] (c0) at (6.50,-0.35) {长度 0};
\draw[flow] (c5) -- node[above, font=\scriptsize] {\(\pi[4]\)} (c3);
\draw[flow] (c3) -- node[above, font=\scriptsize] {\(\pi[2]\)} (c1);
\draw[flow] (c1) -- node[above, font=\scriptsize] {\(\pi[0]\)} (c0);
\end{tikzpicture}
\end{center}
}

\newcommand{\dpchart}{%
\begin{center}
\begin{tikzpicture}[
  mask/.style={draw=green!45!black, fill=green!7, rounded corners=2pt,
    minimum width=1.0cm, minimum height=0.48cm, font=\small},
  flow/.style={-{Latex[length=2mm]}, thick, draw=black!60},
  rulebox/.style={draw=purple!55!black, fill=purple!6, rounded corners=3pt,
    align=left, inner sep=5pt, font=\small}
]
\node[anchor=west, font=\small\bfseries] at (0,0.85) {单行可用掩码};
\node[mask] (m0) at (0,0.15) {\texttt{000}};
\node[mask] (m1) at (1.18,0.15) {\texttt{001}};
\node[mask] (m2) at (2.36,0.15) {\texttt{010}};
\node[mask] (m3) at (3.54,0.15) {\texttt{100}};
\node[mask] (m4) at (4.72,0.15) {\texttt{101}};
\node[rulebox] (rule) at (9.00,0.15) {\textbf{相邻两行的兼容条件}\\[2pt]
  \(s \mathbin{\&} t = 0\)：保留\qquad
  \(s \mathbin{\&} t \ne 0\)：跳过};
\draw[flow] (m4.east) -- (rule.west);
\node[anchor=west, font=\scriptsize, text=black!65] at (0,-0.45)
  {同一行先排除相邻的 1；相邻行再检查同一列是否同时取 1。};
\end{tikzpicture}
\end{center}
}

\newcommand{\segchart}{%
\begin{center}
\begin{tikzpicture}[
  seg/.style={draw=red!55!black, fill=red!5, rounded corners=3pt,
    minimum width=2.8cm, minimum height=0.68cm, align=center, font=\small},
  flow/.style={-{Latex[length=2mm]}, thick, draw=black!65},
  note/.style={draw=orange!65!black, fill=orange!7, rounded corners=3pt,
    align=center, inner sep=4pt, font=\small}
]
\node[seg] (root) at (3.20,1.10) {父结点\quad 区间总和\\懒标记 \((\mathit{mul},\mathit{add})\)};
\node[seg] (left) at (1.25,-0.05) {左子结点\\区间总和};
\node[seg] (right) at (5.15,-0.05) {右子结点\\区间总和};
\draw[flow] (root.south west) -- node[above left, font=\scriptsize] {\texttt{apply(k,b)}} (left.north);
\draw[flow] (root.south east) -- node[above right, font=\scriptsize] {\texttt{apply(k,b)}} (right.north);
\node[note] (reset) at (3.20,-1.30) {下传完成后清除父结点标记\\\((\mathit{mul},\mathit{add})\gets(1,0)\)};
\draw[flow] (root.south) -- (reset.north);
\end{tikzpicture}
\end{center}
}

\pagestyle{fancy}
\fancyhf{}
\fancyfoot[C]{\xeCJKsetup{CJKecglue={\hskip0pt}}{\songti 折工 C5P-S 2026 参考答案与解析}\\[2pt]{\songti 第\thepage 页，共\pageref{LastPage}页}}
\renewcommand{\headrulewidth}{0pt}
\renewcommand{\footrulewidth}{0pt}

\begin{document}

\begin{center}
{\dengxian
{\fontsize{15.95}{1}\selectfont 2026 €€£超专业级别硬件能力认证第一轮}\\[1em]
{\fontsize{14.05}{1}\selectfont （C5P-S1）提高级 C++语言试题 \textbf{参考答案与解析}}\\[1em]
{\fontsize{12}{1}\selectfont 认证时间：2026年9月18日14:30\textasciitilde16:30}\\[1em]
{\fontsize{12}{1}\selectfont 编者：唐一潇、季灵、裘泽燑}
}
\end{center}

\vspace{6pt}
\noindent\textbf{考生注意事项：}
\begin{itemize}
  \item 试题纸共有 16 页，答题纸共有 1 页，满分 100 分。请在答题纸上作答，写在试题纸上的一律无效。
  \item 不得使用任何电子设备（如计算器、手机、电子词典等）或查阅任何书籍资料。
\end{itemize}

\noindent\textbf{说明：}本文件对应自制模拟卷，不代表任何官方机构。题面、代码、解析和考点矩阵按本套 42 题版本同步。

\section*{分值与题号}
\begin{center}
\begin{tabular}{c c c c}
\toprule
部分 & 题号 & 题量与计分 & 小计\\
\midrule
单项选择 & 1--15 & $15\times2$ 分 & 30 分\\
阅读程序（1）KMP & 16--20 & $3\times1.5+1\times3+1\times4$ 分 & 11.5 分\\
阅读程序（2）状态压缩 DP & 21--26 & $2\times1.5+1\times2+3\times3$ 分 & 14 分\\
阅读程序（3）线段树 & 27--32 & $3\times1.5+2\times3+1\times4$ 分 & 14.5 分\\
完善程序 & 33--42 & $10\times3$ 分 & 30 分\\
\midrule
\multicolumn{3}{r}{总分} & 100 分\\
\bottomrule
\end{tabular}
\end{center}

\section*{一、单项选择题（1--15）}
\ans{1.}{B}\expn{\texttt{cat} 用于连接并显示文件内容。}
\ans{2.}{C}\expn{含 $m$ 个叶子的二叉树共有 $2m-1$ 个结点，指针域总数为 $4m-2$；非空指针为 $2m-2$，空指针域为 $2m$。}
\ans{3.}{B}\expn{完全 5 叉树高度为 $h$ 时，叶子数为 $\Theta(5^h)$；每个叶子向根访问经过 $\Theta(h)$ 个结点，因此总访问次数为 $\Theta(h5^h)$。}
\ans{4.}{D}\expn{这是长度为 10 的合法括号序列计数，答案为卡特兰数 $C_5=42$。}
\ans{5.}{A}\expn{只有在 B 为空时把 A 整体转移到 B，才能保证最早入队元素位于 B 栈顶。}
\ans{6.}{C}\expn{将第 $i$ 个被选数减去 $i-1$，问题转化为从 15 个位置中任选 6 个，答案为 $\binom{15}{6}=5005$。}
\ans{7.}{B}\expn{\texttt{vector} 在头部插入或删除需要移动元素，通常为 $O(n)$；\texttt{deque} 支持常数时间随机访问，且两端插入、删除为 $O(1)$；\texttt{set} 不允许重复键；\texttt{priority\_queue::top()} 只读取堆顶，不会删除元素。}
\ans{8.}{C}\expn{拓扑序要求每条有向边的起点排在终点之前；有环图不存在拓扑序。}
\ans{9.}{C}\expn{模式串 \texttt{ababa} 在文本串中从下标 0、2、4、6 开始匹配，允许重叠，起始位置之和为 $0+2+4+6=12$。}
\ans{10.}{B}\expn{SPFA 的最坏时间复杂度可达到 $O(nm)$；Dijkstra 要求边权非负，邻接表配二叉堆时复杂度为 $O((n+m)\log n)$；Bellman--Ford 可以判断从源点可达的负环；Floyd--Warshall 的时间复杂度为 $O(n^3)$，因此只有 B 正确。}
\ans{11.}{D}\expn{模质数 $p$ 下，$1/k+1/(p-k)\equiv0$。括号内缺少 $1$，因此为 $-1$；指数 $p$ 为奇数，结果仍为 $p-1=20260912$。}
\ans{12.}{B}\expn{递归式中 \(a=7\)、\(b=3\)，所以 \(n^{\log_b a}=n^{\log_3 7}\)。非递归项 \(5n\log_2 n+3n=\Theta(n\log n)\)，而 \(\log_3 7\approx1.771\)，它比 \(n^{\log_3 7}\) 小一个多项式量级，属于主定理第一种情况。因此 \(T(n)=\Theta(n^{\log_3 7})\)。}
\ans{13.}{B}\expn{枚举 9 条边中删除 2 条，保留至少一条从 1 到 8 的路径，可行方案共有 15 种。}
\ans{14.}{B}\expn{用容斥计算：$200-100-66-28+33+14+9-4=58$。}
\ans{15.}{A}\expn{由 $A(3,n)=2^{n+3}-3$，可得 $A(3,3)=61$。又因 $2^{50}\equiv-1\pmod {101}$，且 $10^{18}+3\equiv3\pmod {100}$，所以 $A(3,10^{18})\equiv2^3-3=5\pmod {101}$，答案为 \texttt{(61,5)}。}

\section*{二、阅读程序（16--32）}

\subsection*{（1）KMP 前缀函数}
\overview{程序先求每个前缀的最长真前后缀长度，再沿着前缀函数回退，累加整个字符串及各级相同前后缀的长度。}
\kmpchart
\ans{16.}{A（正确）}\expn{程序第 10--16 行计算前缀函数。输入为 \texttt{ababa} 时，各位置的前缀函数值为 \texttt{0,0,1,2,3}，所以第 16 行得到 \texttt{pi[4]=3}。}
\ans{17.}{C}\expn{输入 \texttt{ababa} 时，前缀函数回退链为长度 $5\to3\to1$，第 19--20 行累加得到 $5+3+1=9$。}
\ans{18.}{A}\expn{输入 \texttt{bbba} 时，按第 12 行的 \texttt{while} 计算得到输出 4；改为只回退一次的 \texttt{if} 后，第 12 行得到错误的前缀函数，输出为 5，因此二元组为 \texttt{(4,5)}。}
\ans{19.}{C}\expn{程序第 19--20 行沿前缀函数回退链累加。该字符串得到的相同前后缀长度为 $14,8,2$，其和为 $24$，由第 21 行输出。}
\ans{20.}{B}\expn{程序第 10--16 行计算前缀函数时，\texttt{j} 的增加和回退总次数都是线性级别；第 19--20 行沿回退链访问的长度也不超过 $n$，所以时间复杂度为 $O(n)$。第 3--5 行的数组和字符串占用 $O(n)$ 空间。}

\subsection*{（2）状态压缩 DP}
\overview{每一行用无相邻 1 的位掩码表示选取位置，转移要求相邻两行掩码不相交。}
\dpchart
\ans{21.}{A（正确）}\expn{第 16--18 行保留没有相邻 1 的状态；三列中合法状态为 \texttt{000,001,010,100,101}，共 5 个。}
\ans{22.}{A（正确）}\expn{第 30--31 行判断两行是否冲突。\texttt{101\&001=001} 非零，所以执行第 31 行的 \texttt{continue}，不会执行后面的转移累加。}
\ans{23.}{A}\expn{输入为三行两列的全空棋盘时，保留第 26 行的输出为 17；删除第 26 行后，上一轮滚动数组中的数会残留，输出变为 23。}
\ans{24.}{B}\expn{程序第 16--18 行预处理状态，复杂度为 $O(2^m)$；第 26--34 行对每一行枚举状态对，复杂度为 $O(nV^2)$，总复杂度为 $O(nV^2+2^m)$。}
\ans{25.}{C}\expn{程序第 26--34 行进行状态转移。两行三列全可选时，按行状态转移统计得到 17。}
\ans{26.}{A}\expn{输入为三行两列的全空棋盘时，保留第 23--24 行可正确交换滚动数组，输出为 17；删除这两行后第一轮清空了含初值的数组，后续转移全为 0，输出为 0。}

\subsection*{（3）区间仿射更新线段树}
\overview{每个结点保存所代表区间的总和。一次完整覆盖的更新可以先记在当前结点，之后访问子区间时再传给两个子结点。}
\segchart
\ans{27.}{A（正确）}\expn{第 38--40 行是 \texttt{upd} 的完整覆盖分支：第 39 行调用一次 \texttt{apply}，第 40 行立即返回，不再访问子结点。}
\ans{28.}{B（错误）}\expn{第 32、33 行下传懒标记，第 34、35 行分别重置 \texttt{mul} 和 \texttt{add}。删除第 35 行会使旧的加法标记重复下传，导致结果错误。}
\ans{29.}{A（正确）}\expn{程序第 12--22 行建树为 $O(n)$，第 37--49 行的每次更新和第 50--61 行的每次查询均为 $O(\log n)$，总量可写成 $O(n+q\log n)$，也属于 $O((n+q)\log n)$。}
\ans{30.}{B}\expn{程序第 39 行调用 \texttt{apply} 更新区间总和。区间中共有 \texttt{r-l+1} 个数，每个数先乘 \texttt{k} 再加 \texttt{b}，所以总和变为 \texttt{tr[p].sum*k+(r-l+1)*b}，最后取模。}
\ans{31.}{A}\expn{第 32、33 行分别对左右子结点调用 \texttt{apply}。交换调用顺序只改变访问顺序，不改变结果；其余修改会丢失或重复应用更新。}
\ans{32.}{B（4 分）}\expn{三次更新后的数组为 \texttt{8 11 27 39 13 13}，因此区间 $[2,6]$ 的和为 $11+27+39+13+13=103$。}

\section*{三、完善程序（33--42）}
\subsection*{（1）三色序列}
\ans{33.}{C}\expn{先把前一个位置的计数复制过来；如果当前字符属于颜色 \texttt{t}，再多计数一次，因此填 \texttt{pre[t][i-1] + (t == c)}。}
\ans{34.}{B}\expn{\texttt{f} 记录的是交换次数，初始状态尚未构造任何字符，代价应为 0。}
\ans{35.}{A}\expn{前 \texttt{len} 个位置已经构造，其中使用了 \texttt{r} 个 R 和 \texttt{g} 个 G，所以 Y 的使用数为 \texttt{len-r-g}。}
\ans{36.}{A}\expn{新字符不能和上一字符相同，也不能超过该颜色的总数；满足任一条件就跳过，因此填 \texttt{c == last || used[c] == cnt[c]}。}
\ans{37.}{C}\expn{当前状态代价为 \texttt{val}，加入该字符的增量为 \texttt{w}，所以新值为 \texttt{val + w}。}

\subsection*{（2）盒子}
\ans{38.}{B}\expn{\texttt{lo} 保存当前选中的较小间隙；若数量过多，应移出其中最大的一个，因此填 \texttt{prev(lo.end())}。}
\ans{39.}{B}\expn{\texttt{x} 是 \texttt{lo} 中最大的间隙，\texttt{y} 是 \texttt{hi} 中最小的间隙。只有在 \texttt{*x <= *y} 时才不需要交换。}
\ans{40.}{B}\expn{交换后，\texttt{lo} 中的 \texttt{vx} 被 \texttt{vy} 替换，答案变化为 \texttt{vy-vx}；代码用 \texttt{ans +=} 累加这个变化量。}
\ans{41.}{A}\expn{\texttt{n} 个盒子最终形成 \texttt{k} 个链，每条链需要相邻盒子之间的间隙，共需保留 \texttt{n-k} 个间隙。}
\ans{42.}{C}\expn{删除体积为 \texttt{v} 的中间盒子后，左、右盒子直接相邻，新间隙为 \texttt{r-l}。}

\end{document}
